\documentclass[twocolumn,a4paper,aps,amsmath,amssymb,floatfix,superscriptaddress]{revtex4-2}
\usepackage[utf8]{inputenc}
\usepackage{booktabs}
\usepackage{graphicx}
\usepackage{amsmath, amssymb, amsthm}
\usepackage{physics}
\usepackage[shortlabels]{enumitem}
\usepackage{cancel}
\usepackage{lastpage}
\usepackage{braket}
\usepackage{mathrsfs}
\usepackage{appendix}
\usepackage{listings}
\usepackage{xcolor}
\usepackage{graphicx}
\usepackage{dcolumn}
\usepackage[hidelinks]{hyperref}
\usepackage{bm}
\usepackage[utf8]{inputenc}

\definecolor{codegreen}{rgb}{0,0.6,0}
\definecolor{codegray}{rgb}{0.5,0.5,0.5}
\definecolor{codepurple}{rgb}{0.58,0,0.82}
\definecolor{backcolour}{rgb}{0.95,0.95,0.92}

\newcommand{\todo}[1]{
    \textbf{TODO: $<$#1$>$}
}

\usepackage{fancyhdr}
\pagestyle{fancy}
\lhead{NS-204B - Statystische fysica theorie en experiment}
\rhead{}
\fancyfoot{}
\fancyfoot[R]{Page \thepage /\pageref{LastPage}}

\newcommand{\R}{\mathbb{R}}
\newcommand{\p}{\mathbb{P}}
\newcommand{\N}{\mathbb{N}}
\newcommand{\Z}{\mathbb{Z}}
\newcommand{\Q}{\mathbb{Q}}
\newcommand{\I}{\mathbb{I}}
\newcommand{\E}{\mathbb{E}}

\newcommand{\mbeq}{\overset{!}{=}}
\newcommand{\ko}{\frac{1}{4\pi \epsilon_0}}
\newcommand{\fp}{../figures/}

\lstset{frame=tb,
  language=java,
  aboveskip=3mm,
  belowskip=3mm,
  showstringspaces=false,
  columns=flexible,
  basicstyle={\small\ttfamily},
  numbers=none,
  numberstyle=\tiny\color{gray},
  keywordstyle=\color{blue},
  commentstyle=\color{dkgreen},
  stringstyle=\color{mauve},
  breaklines=true,
  breakatwhitespace=true,
  tabsize=3
}

\lstdefinestyle{mystyle}{
    backgroundcolor=\color{backcolour},   
    commentstyle=\color{codegreen},
    keywordstyle=\color{blue},
    numberstyle=\tiny\color{codegray},
    stringstyle=\color{codegreen},
    basicstyle={\small\ttfamily},,
    breakatwhitespace=false,         
    breaklines=true,                 
    captionpos=b,                    
    keepspaces=true,                 
    numbers=left,                    
    numbersep=5pt,                  
    showspaces=false,                
    showstringspaces=false,
    showtabs=true,                  
    tabsize=1
}

\lstset{style=mystyle}

\graphicspath {{../../figures/}}

\begin{document}

\title{They Like to Move It, Move It  — \\ Determining the Viscosity of Colloidal Suspensions of Monodisperse Polystyrene Microspheres in Water Using Brownian Motion}
        
\author{Bela F. Brunner}
%Lines break automatically or can be forced with \\
\affiliation{Department of Physics, \\
Utrecht University}

\author{Adam Zich}
\affiliation{Department of Physics, \\
Utrecht University}

\date{\today}

\begin{abstract}
    Abstract
\end{abstract}

\maketitle

\section{Introduction}\label{intro}
    Brownian motion describes the continuous, random movement of microscopic particles suspended in a fluid. The phenomenon was first observed by Robert Brown in the nineteenth century, who noticed the irregular movement of pollen grains suspended in water. Its theoretical foundations were subsequently developed by Albert Einstein providing a physical explanation for the observed motion and establishing its connection to the microscopic properties of fluids \cite{blundell2009}.

    Einstein's theoretical description relates the mean squared displacement of suspended particles to the diffusion coefficient and elapsed time \cite{einstein1905}. This relationship provides a way to investigate the properties of a fluid through the motion of particles suspended within it. In particular, the Stokes-Einstein relation relates the diffusion coefficient to the particle size, fluid viscosity and temperature \cite{britannica_brownian_motion}. Consequently, by tracking the Brownian motion of particles with known properties, the viscosity of the surrounding fluid can be determined.
    
    In this experiment, Brownian motion is used to determine the viscosity of aqueous suspensions containing monodisperse polystyrene microspheres of different sizes. By tracking the displacement of individual particles over time, their diffusion coefficients can be determined and related to the viscosity of the surrounding fluid. The main objective is to investigate whether this method provides a reliable determination of viscosity and to examine how the measured viscosity depends on the size of the suspended microspheres. If successfull, this would allow an \textit{in situ}\footnote{Meaning without altering the position of the liquid} measurement of the viscosity of the liquid. This is especially advantageous since it allows for measurements without disturbing the microsystem \cite{kimBrownianMicroscopySimultaneous2008}.
    
    


\section{Theory}\label{theory}
    A key relation connecting Brownian motion to the fluid's viscosity is the Stokes–Einstein relation, which is given as follows
    \begin{align*}
        D=\frac{k_BT}{6\pi \eta r}
    \end{align*}
    where $D$ in $\text{m}^2/\text{s}$ is the diffusion coefficient, $k_B$ is the Boltzmann constant given as $1.38 \times 10^{-23} \, \text{J/K}$, $T$ is the thermodynamic temperature in Kelvin, $\eta$ is the dynamic viscosity in Pascal-seconds, and $r$ is the particle radius in meters \cite{einstein1905}. It applies to spherical particles suspended in a Newtonian fluid and provides us with a way to determine the viscosity of the fluid by measuring the diffusion coefficient of the suspended particles. These conditions are mostly fufilled by our setup. We can rewrite this equation in terms of the viscosity as
    \begin{align*}
        \eta = \frac{k_BT}{6\pi Dr}.
    \end{align*}
    To find this, we can use the relation between the mean squared displacement of a particle and the diffusion coefficient. 
    
    In Brownian motion we assume that the motion of each distinct particle is independent of the motion of the other suspended particles. If we then also impose the condition that the number of particles per unit volume $\nu$ is a function of only $x$ and $t$, i.e. $\nu = f(x, t)$, we get that it sattisfies the equation
    \begin{equation}
        \frac{\partial f}{\partial t} = D\frac{\partial^2 f}{\partial x^2},
    \end{equation}
    where $D$ is the diffusion coefficient \cite{einstein1905}.\\
    We then define $f(x, t)dx$ to denote the number of particles that have moved by a displacement between $x$ and $x + dx$ in the $x$ direction in the time interval $[0, t]$. We can impose the initial conditions as
    \begin{align*}
        f(x, 0) = 0,\\
        \int_{-\infty}^\infty f(x, 0)dx = n,
    \end{align*}
    where $n$ is the total number of particles suspended in the liquid \cite{einstein1905}. The solution to $f(x, t)$ can in this case be found analytically and is given by
    \begin{equation}
        f(x, t) = \frac{n}{\sqrt{4\pi Dt}}e^{-\frac{x^2}{4Dt}}.
    \end{equation}
    We can then determine the average displacement of a particle along the $x$ direction as
    \begin{equation}
        \lambda_x = \sqrt{\braket{x^2}} = \sqrt{2Dt},
    \end{equation}
    where $\braket{x^2}$ is the arithmetic mean of squares of displacements along the $x$ direction \cite{einstein1905}.\\
    In the 2D case, we take the effect to be independent in the $x$ and $y$ axes. This then gives us an expression for the mean squared displacement in the $xy$-plane as
    \begin{equation}
        \braket{R^2} = \braket{x^2} + \braket{y^2} = \braket{x^2 + y^2} = 2Dt + 2Dt = 4Dt.
    \end{equation}
    We can then use this relation to find an equation for the diffusion coefficient as
    \begin{align}
        4Dt &= \braket{R^2},\\
        \frac{d}{dt}4Dt &= \frac{d}{dt}\braket{R^2},\\
        4D &= \frac{d}{dt}\braket{R^2},\\
        D &= \frac{d\braket{R^2}}{4dt}. \label{diff_coef}
    \end{align}
    We can then get an expression for the dynamic viscosity as
    \begin{equation}
        \eta = \frac{k_BT}{6\pi Dr} = \frac{2k_BT}{3\pi \frac{d\braket{R^2}}{dt} r}. \label{viscosity}
    \end{equation}
    Calculating the uncertainty of a quantity is done using the formula
    \begin{equation}
        \sigma_A = A \cdot \sqrt{\sum_n \sigma_n^2}, \label{uncertainty}
    \end{equation}
    where $A$ represents the quantity we want to calculate the resulting ucnertainty for and $\sigma_n$ are relative uncertainties of quantities that are in the formula for $A$.

\section{Methodology}
    As samples to be analysed, we used monodisperse polystyrene microspheres with factory-specified diameters of $0.75-2 \mu\text{m}$, suspended in water \cite{PolybeadMicrospheres200um}. The resulting liquid is a colloidal suspension, which describes mixtures in which the colloidal particles are significantly larger than the solvent molecules but small enough to exhibit Brownian motion \cite{dijkstraComputerSimulationsCharge2001}. In our case, the particles are suspended at a concentration significantly lower than that of the stock suspension in which they are supplied, namely a $2.5\%$ solids (w/v) aqueous suspension \cite{PolybeadMicrospheres200um}. We used samples with a mass-per-volume concentration of $0.1\%$ across all measurements to minimise the effects of concentration as an additional variable and to keep the results comparable. These samples were prepared by diluting the stock suspension with water.

    We prepared samples with four different particle radii, namely $r \in \{0.373, 0.53, 0.77, 1.03\} \mu m$, for which the differences from the factory-specified values had been determined in advance. The samples were then placed on a microscope slide and covered with a coverslip to prevent evaporation of the liquid. The slide was placed under a microscope with a magnification of $40\times$, and a light source was positioned at an angle to the sample to create contrast and make the particles more visible.

    To calculate the final viscosity using Eqn. \ref{viscosity}, we needed to determine the temperature and diffusion coefficient. The temperature was measured using a room thermometer located in the room where the experiment was conducted. This thermometer had a precision of $0.1^\circ C$, and the room temperature ranged between $20.9^\circ C$ and $21^\circ C$.

    We calculated the diffusion coefficient by measuring the mean squared displacement $\braket{R^2}$ of the particles over time. To do this, we recorded videos of the particles using the DinoCapture 2.0 software. The videos were recorded at a frame rate of 20 frames per second, with durations between 3 and 9 seconds and a resolution of 640x480 pixels. The resulting videos were then processed using the Python library Trackpy to track the movements of the particles and calculate their mean squared displacement over time.

    For the tracking to work correctly, several parameters had to be tailored to each video. One required parameter was an estimate of the particle size in pixels, called the "feature size". This value was mostly kept uniform across videos of particles with the same diameter. In some cases, however, the microscope focus required the value to be adjusted for Trackpy to track the particles correctly. We identified this as an apparent diffraction effect, in which a bright spot appeared in the middle of the particle due to interference. This meant that a higher feature size was needed to track the entire particle and avoid tracking only part of it. Another important parameter was the so-called "minmass". This parameter determines how strong the particle signal must be to be detected and tracked. By setting a higher minmass, particles had to be more clearly visible and in the foreground to be tracked. This helped eliminate accidental tracking of shadows from out-of-focus particles, as well as dust present in the setup. To identify these particles, we visually inspected the tracking results and adjusted the parameters until the tracking was clear throughout the video. We also implemented drift correction to eliminate systematic particle movement caused by external factors, such as vibrations or air currents. This was done by calculating the average displacement of all tracked particles and subtracting it from the individual particle displacements.

    To convert the tracked mean squared displacement from pixels to micrometres, we needed to determine a calibration factor expressing how many micrometres corresponded to one pixel. For this, we used the software Fiji and a micrometer ruler placed under the microscope with the same magnification and camera settings as those used for our videos. We then took a picture and measured the distance between eight marks on the ruler in pixels, while the corresponding physical distance was known. We estimated the uncertainty from the width of one mark, which yielded a calibration factor of $0.130 \pm 0.003 \mu m/px$.

    The mean squared displacement (MSD) was calculated from the tracked particle trajectories using the \texttt{emsd} function from the Python library Trackpy. This function calculates the ensemble mean squared displacement by averaging squared particle displacements over the available trajectories and time origins for each time lag. The pixel-to-micrometre calibration factor and the video frame rate were supplied to convert the resulting displacements into physical units and time. The maximum time lag was set to one quarter of the total number of video frames to limit the analysis to time intervals with sufficient tracking data. The resulting MSD values were then used to determine the diffusion coefficient from a linear fit.

    To obtain the diffusion coefficient from the mean squared displacement, we used the formula in Eqn. \ref{diff_coef}. We fitted a linear function to the mean squared displacement as a function of time to obtain a slope (derivative) equal to $4D$. We constrained the fit to pass through the origin, in accordance with the theoretical expectation that the MSD is zero at zero time lag. The diffusion coefficient was then obtained by dividing the fitted slope by four. This constraint assumes that any experimental offset in the MSD is negligible; its possible influence on the fitted diffusion coefficient should therefore be considered when interpreting the results.

    The uncertainties propagated through this process included the uncertainty in the particle diameter, the uncertainty arising from the thermometer precision, the uncertainty in the micrometre-to-pixel conversion, and the uncertainty in the diffusion coefficient obtained from the linear fit. The uncertainty in the viscosity, as well as all preceding uncertainty propagations, was calculated using the formula in Eqn. \ref{uncertainty}.

    Once the diffusion coefficient had been calculated, we used the Stokes-Einstein relation (Eqn. \ref{viscosity}) to determine the resulting viscosity for a given particle radius and its uncertainty. To average the different viscosity values calculated for the same radius, we used the Birge ratio method to account for a possible mismatch between the individual uncertainties and the spread of the viscosity values \cite{trassinelli2026}.

\begin{figure}[htbp]
    \includegraphics[width=0.5\textwidth]{\fp setup_schematic.png}
    \caption{Schematic of the experimental setup used to record the Brownian motion of polystyrene microspheres suspended in water. The sample was placed on a microscope slide under a $40\times$ microscope objective and illuminated from an angle to improve particle contrast. Particle motion was recorded using a camera connected to a computer.}
    \label{fig:setup_schematic}
\end{figure}
```
\section{Results}\label{results}
    Our meausred value of temperature was of the order of $\sim21^\circ C$. The diffusion coefficient was of the order $\sim 10^{-1}\mu m$ depending on the particle radius. In the end we calcualted 25 viscosity values, 6 for $r = 0.373$, 6 for $r = 0.53$, 8 for $r = 0.77$ and 5 for $r = 1.03$. We then averaged the results for each radius. As we did not expect the viscosity to change, we performed a constant function fit through our data and obtained a vlaue of viscosity of $\nu = 2.4 \pm 0.1 mPa\cdot s$ as can be seen in Figure \ref{fig:result}.
    \begin{figure}[htbp]
        \includegraphics[width=0.5\textwidth]{\fp averaged.png}
        \caption{The final result of average measurement of viscosity for particle radii $r \in \{1.03, 0.77, 0.53, 0.373\} \mu m$ and a fit of a cosntant function that yields resulting viscosity of $\nu = 2.4 \pm 0.1 mPa\cdot s$.}
        \label{fig:result}
    \end{figure}

\section{Discussion}
    The aquisition and analysis of data for the experiment outlined in this article proved to be rather difficult.\par
    One of the problems we encountered was the inconsistency of our setup. Since the stand holding the light and the lense were not fixed in a position relative to the microscope, every time we disassembled and rebuilt the setup, our lighting conditions changed. We tried to reduce this effect as much as possible by placing the microscope and light aparatus into a relative position as close to the initial one as we could, however, this process was not aided by any tools and only eye estimation was used. Unfortunately, we were also not able to do perform this in an efficient way at the beginning of the data aquisition phase, which cost us valubale time. A way to achieve a more consistant setup would be using a microscope with a fixed light or potentially securing the light to the microscope in some way.\par
    Another issue, which we think had a bigger contribution to decreasing the quality of final data, was the tracking of particles using trackpy. Even after trying to mitigate effects disturbing the tracking of the particles, we observed that the effect of wrongful tracking is significant, which we were not able to quantify within the uncertainties. Furthermore, this might be a systematic error.\par
    One of the issues with the tracking of these particles lied in the light interfenece effects that made it so the particles transitioned from dark to light around the focus point. We tried to go around this problem by implementing a slight de-focus for the marticles we wanted to track to elimiate this problem, but then the particles were naturally less sharp, which introduced a new problem. Even though we did not quantify the difference beween in- and out of focus particles, we believe the de-focused particles produced better data, or were the very least easier to fine-tune the trackpy parameters for.\par
    Another contribution to the difficulty of tracking was the fact that our microscope was set to a specific depth of field, where the particloes we wanted to observe we in the focus we wanted (which as mentioned before was a slight de-focus). However, the sample was not 2D, but the particles also moved in the direction of our line of sight. This meant that sometimes the particles moved further into de-focusing themselves or covnersely, moved into the focus, introducing the interference problem again. In practice this meant that we had to only keep video data that did not contain these issues, which meant taking more videos than needed and filtering those that were deemed unsuable due to these phenomena.\par
    Finally, another obstacle we had to overcome was the fact that the longer we observed a given sample, the higher the calcualted viscosity seemed to climb. We attributed this to the fact that are sample was a liquid, whcih was only trapped between the sample glass plate on either side. This mean, however, that parts of the liquid could evaporate and other processes disturbing the integrity of the sample could take place. To mitigate this phenomen and its effect on our data, we only used a given sample for a limited time after its preperation, wafter which we were forced to make a new one if more measurements were required.

\section{Conclusions}\label{conclusions}
    Conclusion

\bibliographystyle{unsrt}
\addcontentsline{toc}{section}{Bibliography}
\bibliography{zotero_bibliography.bib}

\end{document}